% Beamer skeleton for pedagogical info decks (see DECK_PROTOCOL.md).
% Rename colors/chain segments per project; keep the structure.
\documentclass[aspectratio=169,10pt]{beamer}

\usetheme{default}
\usecolortheme{seagull}
\setbeamertemplate{navigation symbols}{}
\setbeamertemplate{itemize items}{\textbullet}
\usepackage{amsmath, amssymb, bm, graphicx}

% deck-wide color code: two physical branches + one accent (instability/controls)
\definecolor{cL}{HTML}{1F77B4}   % branch A (e.g. superluminal / L sector)
\definecolor{cA}{HTML}{D62728}   % branch B (e.g. subluminal / A sector)
\definecolor{cX}{HTML}{FF7F0E}   % accent: instability / positive control
\definecolor{cInk}{HTML}{202030}
\setbeamercolor{frametitle}{fg=cInk}
\setbeamerfont{frametitle}{series=\bfseries,size=\large}
\setbeamercolor{structure}{fg=cL}

% footline roadmap: the deck's spine with the current segment bolded.
% Define one \chainN command per lecture; call it at each lecture boundary.
\newcommand{\chainstrip}{}
\setbeamertemplate{footline}{%
\hspace{3mm}{\tiny\chainstrip}\hfill {\tiny\insertframenumber\,/\,\inserttotalframenumber}\hspace{3mm}\vspace{1.2mm}} \newcommand{\chainI}{\renewcommand{\chainstrip}{\textbf{step one} \textcolor{gray}{$\to$ step two $\to$ step three $\to$ step four}}} \newcommand{\chainII}{\renewcommand{\chainstrip}{\textcolor{gray}{step one $\to$} \textbf{step two} \textcolor{gray}{$\to$ step three $\to$ step four}}}
% ... one per lecture

% the two content boxes of the protocol
\newcommand{\analog}[1]{\vfill\begin{beamercolorbox}[rounded=true,sep=4pt]{block body} {\footnotesize \textbf{Analog.} #1}\end{beamercolorbox}} \newcommand{\defbox}[1]{\begin{beamercolorbox}[rounded=true,sep=4pt]{block body} {\footnotesize \textbf{Definition.} #1}\end{beamercolorbox}}

% project math macros go here (match the source manuscript's notation)

\title{...}   % title only: no subtitle, no self-description (DECK_PROTOCOL.md §2)
\author{}
\date{}

% Frame patterns:
%
% Assertion-titled content frame:
%   \begin{frame}{The claim of this frame, as a sentence}
%   [context sentence] [derivation with intermediate line] [physical reading]
%   \begin{itemize}...\end{itemize}
%   \analog{named system: the mapping.}
%   \end{frame}
%
% Tool frame (method introduction, before first use):
%   \begin{frame}{Tool: what a ... is}
%   \defbox{...definition with every symbol named...}
%   [how it is instantiated here]
%   \end{frame}
%
% Figure frame with side-column reading guide:
%   \begin{frame}{Assertion about what the figure shows}
%   \begin{columns}[T]
%   \column{0.30\textwidth} {\footnotesize \textbf{Left:} ... \textbf{Right:} ...}
%   \column{0.68\textwidth} \includegraphics[height=0.86\textheight]{plot_deck_sN}
%   \end{columns}
%   \end{frame}
%
% Lecture recap frame:
%   \begin{frame}{Lecture N in one line} \centering \Large ... \end{frame}
